\scanpage{63}
$D,D',\Delta,\Delta'$ est en position générale, donc en
position \unsure{standard|standart}) que $\overrightarrow{tu\unsure{v|u}}$ et $\overrightarrow{tu'v'}$ se correspondent.
[Plus généralement, si on a $k$ droites
$\Delta_1\ldots\Delta_k$ passant par $s$ et non par $t$, \inserted{et coupant $D$ en $u_i$ et $D'$ en $u'_i$,} et si
on prend sur $D\setminus\{t\}$ et $D'\setminus\{t'\}$ les relations
d'ordre \inserted{induites par des orientations de $D,D'$} qui se correspondent, alors si on ordonne
les $\Delta_i$ de façon que les \deleted{$D$} $u_i$ aillent en
croissant, de \unsure{m\textsuperscript{e}|$=$} les $u'_i$.] Ceci explicite la
bijection $\Omega(D)\simeq\Omega(D')$ en termes de la
combinatoire du $K^1$ et de ses ps-droites
sauf dans le cas où il n'y a qu'une seule ps-droite \inserted{$\Delta$} passant par $s$ et non par $t$, (ce qui implique
qu'il existe \unsure{exact$^{\mathrm{mt}}$|exactement} 2 psdr passant par $s$, i.e. que
l'ordre de $s$ est 4 \unsure{\}|)}, et qu'il existe une ps.dr. $\Delta'$
passant par $s$ et par $t$) Supposons qu'il existe une
ps.dr. $D''$ \inserted{ne passant pas par $t$} distincte des \unsure{$\Delta'$ en|$\Delta'$} (i.e. que la figure
% slanted
\margin{NB $D''$ ne passe pas par $s$ i.e. cette ps.dr. $\ne\Delta,\Delta'$}
ne se réduise à un faisceau des ps.dr. passant
par $t$, et une sécante $\Delta$ de ce faisceau) --
elle \deleted{est} \deleted{\unsure{passe}} \deleted{\unsure{donc}} $\ne\Delta,\Delta',D,D'$ et ne passe pas par $s$
Soient $x,x'$ les pts d'int. de $D''$ avec $D,D'$ et $u,u'$
ceux de $\Delta$ avec $D,D'$. Si on a $x\ne u$, $x'\ne u'$, alors
on peut comparer par ce qui précède les \deleted{\unsure{orientations}} \inserted{ensembles}
% archive sheet number: 62
